\documentclass[11pt]{article}
\newcommand{\runningtitle}{Signed zero multiplicities and Gaussian prime angles}
\usepackage[T1]{fontenc}
\usepackage{lmodern,amsmath,amsthm,amssymb,booktabs,microtype,xcolor}
\usepackage[letterpaper,margin=.94in,top=.83in,bottom=.86in]{geometry}
\usepackage[colorlinks=true,linkcolor=black,urlcolor=teal!65!black,citecolor=black,breaklinks=true]{hyperref}
\usepackage{fancyhdr}
\pagestyle{fancy}
\fancyhf{}\lhead{\small\sffamily\runningtitle}\rhead{\small\sffamily Research note}\cfoot{\small\thepage}
\setlength{\headheight}{14pt}
\setlength{\emergencystretch}{3em}
\newtheorem{theorem}{Theorem}
\newtheorem{lemma}{Lemma}
\newcommand{\paperdate}{November 25, 2026}
\date{\paperdate}
\author{}
\raggedbottom

\hypersetup{pdftitle={Signed multiplicities in short intervals and optimal weights for Gaussian prime angles}}
\title{Signed multiplicities in short intervals and optimal weights for Gaussian prime angles}
\begin{document}
\maketitle
\begin{abstract}
The one-sided prime-sum criterion for Dirichlet GRH records differences of zero orders. We use this observation to compare zero multiplicities of arbitrary fixed finite families in every sufficiently high interval of length $(2/\pi+\epsilon)\log\log\log T$. A second result places both signs of the twist-family difference in any prescribed box of prime phases at polynomial interval lengths. For Hecke characters of $\mathbb Q(i)$ we solve an extremal problem for positive angular weights and identify the additional prime estimate that would give their Riemann hypotheses.
\end{abstract}

\paragraph{From the formal proof to the two analytic questions.}
The readable \href{https://github.com/lsernau/grh-lean/tree/72fead102f663d5c7a407964a13dc0cf9d977207/grh}{Lean source} builds two estimates for the same lattice of automorphic forms. Short vectors with a spanning estimate bound its covolume from above. An integral determinant bounds it from below. Their comparison in \texttt{star\_inequality} yields a one-sided bound for a smoothed sum over prime powers in the class $1\bmod q$, with a factor $1-(\cdot/d)$. The logarithmic derivatives of Dirichlet functions and their quadratic twists give its transform. At $\rho$ the numerator of the residue is the sum of the untwisted orders minus the sum of the twisted orders. The denominator is $\varphi(q)\rho(1-\rho)$. By choosing $d$ to avoid twisted zeros on a given compact rectangle to the right of $1/2$, the proof obtains the signs needed for Landau's theorem. Section~\ref{sec:short} instead studies the differences on the critical line.

The five public statements in \texttt{Main.lean} concern the Riemann zeta function and Dirichlet characters. The readable release records Lean 4.35.0 and Mathlib commit \texttt{56a00585}, with the standard axioms \texttt{propext}, \texttt{Classical.choice} and \texttt{Quot.sound}. We inspected the pinned sources and used the published checking record. The proofs below are on paper. Their finite checks are described in the archive.
\section{Multiplicity differences in intervals of triple-logarithmic length}
\label{sec:short}
For distinct primitive Dirichlet characters $\chi_1,\ldots,\chi_s$ of conductors
$q_1,\ldots,q_s$, let $m_j(t)$ be the order of $L(s',\chi_j)$ at $s'=1/2+it$.
The conductor-one character gives the Riemann zeta function. Fix $s\ge2$.
For $c\in\mathbb R^s$ put $z_c(t)=\sum_jc_jm_j(t)$ and
$z_{c,\pm}=\max(\pm z_c,0)$. Write $\log_3 T=\log\log\log T$.
Sums below count each ordinate once.

\begin{theorem}\label{th:short}
Fix $\epsilon>0$ and the characters. For all sufficiently large $T$, every
\[
 (2/\pi+\epsilon)\log_3 T\le H\le T
\]
satisfies, simultaneously for both signs and every $c$ with
$\sum_jc_j=0$, $\|c\|_2=1$,
\[
 \sum_{T<t\le T+H}z_{c,\pm}(t)\ \gg_{\epsilon,\chi_1,\ldots,\chi_s}\ H\log\log H .
\]
\end{theorem}
For two characters take $c=(1,-1)/\sqrt2$. Each interval contains both an
ordinate where the first multiplicity is greater and one where the second is
greater. The constant $2/\pi$ is sufficient as a limiting coefficient in
this construction. Its optimality is open here.

\paragraph{Comparison.}
Srinivas proved that for distinct functions in the Selberg class of
equal degree, each direction of excess occurs within distance
$O(\log\log T)$ of every large $T$, using a kernel of double-exponential
decay. Our use of real critical-line ordinates gives a smaller interval for
fixed Dirichlet families and a lower bound for the mass in it.
Banks gives bounds uniform in the character modulus.
Fujii and Bombieri--Perelli obtain much larger total multiplicity counts in
global ranges. Banks and Loftus detect genuinely unshared zeros in longer
windows. In particular, if the multiplicities at a common zero are 2 and
1, our sum counts one unit there. The conjectural scarcity of common zeros
remains open. See \href{https://doi.org/10.4064/aa103-3-1}{Srinivas},
\href{https://arxiv.org/abs/2302.07073}{Banks},
\href{https://matwbn.icm.edu.pl/ksiazki/aa/aa28/aa2847.pdf}{Fujii},
\href{https://doi.org/10.4064/aa-83-3-271-281}{Bombieri--Perelli}, and
\href{https://arxiv.org/abs/2607.22930}{Banks--Loftus}.

\subsection{A positive weight and its prime frequencies}
Let $Q$ be the least common multiple of the conductors and
$C(n)=\sum_jc_j\chi_j(n)$. The induced characters of
$(\mathbb Z/Q\mathbb Z)^\times$ are distinct by uniqueness of the primitive
inducer. Character orthogonality gives the mean square of $C(v)$ over the
group as $1$. Choose $v$ with $|C(v)|\ge1$, and choose $\theta$ with
$e^{i\theta}C(v)=|C(v)|$. Put $Y=\lambda\log H$ for a sufficiently
small fixed $\lambda>0$. Take all primes $p\le Y$ congruent to $v$ modulo
$Q$ and write
\[
 w_p=\frac{\log p}{\sqrt p},\qquad V^2=\sum w_p^2
       \sim \frac{(\log Y)^2}{2\varphi(Q)}.
\]
The asymptotic follows by partial summation from the prime number theorem
in the finitely many fixed progressions. Form the real trigonometric sums
\[
 A(t)=\Im\prod_p(1+i(w_p/V)\cos(t\log p+\theta)),\qquad
 P_\pm(t)=1\pm e^{-1/2} A(t).
\]
On the real axis $|A|\le\prod_p\exp(w_p^2/(2V^2))=\sqrt e$, so
$0\le P_\pm\le2$. Extend their finite exponential expansions to complex
$t$. Write any one of these two as $P(t)=\sum_R b_R e^{it\log R}$,
where the positive rationals $R$ are in lowest terms.
Unique factorization gives $b_1=1$ and
$b_p=\pm e^{-1/2}e^{i\theta}w_p/(2V)$ at the chosen primes.
Every other nonconstant frequency, apart from their reciprocals, has at
least two different primes with nonzero exponents and cannot be a prime
power or its reciprocal. Also $b_{1/R}=\overline{b_R}$. Both the numerator
and denominator of $R$ are at most $D=\prod_{p\le Y}p\le e^{2Y}$ for
large $Y$. Expanding the product and using Cauchy--Schwarz gives
$\sum_R|b_R|\le 1+\prod_p(1+w_p/V)\le1+\exp(\sqrt{\pi(Y)})$.

\subsection{An explicit-formula test localized in height}
Choose $0<a<\pi$ and $0<\kappa<1/2$ such that
$a(1/2-\kappa)(2/\pi+\epsilon)>1$. Set $B=\kappa H$, $X=T+H/2$.
Let $k(t)$ be $e^{-\cosh(at)}$ divided by its real integral, and put
\[
 K_B(t)=\frac{1}{2B}\int_{-B}^{B}k(t-y)\,dy,\qquad
 h(t)=P(t)K_B(t-X).
\]
Thus $0\le h(t)\le1/B$ on $\mathbb R$. Use the Fourier convention
$\widehat h(u)=\int_{\mathbb R}h(t)e^{-itu}\,dt$.
For fixed $b$ with $1/2<b<\pi/(2a)$, shifting the contour to $\Im t=\pm b$
gives $|\widehat k(u)|\ll_{a,b}e^{-b|u|}$ since
$|e^{-\cosh(a(t+ib))}|=e^{-\cos(ab)\cosh(at)}$.
Consequently
\[
 \widehat h(u)=\sum_R b_R e^{-iX(u-\log R)}
       \widehat k(u-\log R)\operatorname{sinc}(B(u-\log R)),
 \qquad \operatorname{sinc}(w)=\sin(w)/w,
\]
with value $1$ at zero. Resonant frequencies $u=\log R$ contribute exactly
$b_R$ independently of the center $X$.

For a primitive $\chi_j$ write $a_j=(1-\chi_j(-1))/2$ and
$\psi=\Gamma'/\Gamma$. Under Dirichlet GRH the standard Guinand--Weil
formula in this normalization reads
\begin{align*}
 2\pi\sum_t m_j(t)h(t)
   ={}& 2\pi\mathbf1_{q_j=1}\{h(i/2)+h(-i/2)\}\\
 &+\int_{\mathbb R}h(t)\left[\log(q_j/\pi)+
       \Re\psi\big((1/2+a_j+it)/2\big)\right]dt\\
 &-2\Re\sum_{n\ge2}\frac{\Lambda(n)\chi_j(n)}{\sqrt n}\widehat h(\log n).
\end{align*}
The real-valued test need not be even. The formula follows by integrating the
completed logarithmic derivative on the two sides of the critical strip,
using its functional equation on the left. The strip of decay just proved
justifies both contour shifts and absolute summation on the prime side.
For the explicit-formula convention see
\href{https://doi.org/10.1093/qmath/haab041}{\emph{Fourier optimization and quadratic forms}, section 5}.

After multiplication by $c_j$ and summation, the archimedean contribution
is $O(1)$ uniformly. Indeed, the common $\log(|t|/2)$ from Stirling's
formula cancels since $\sum c_j=0$, the combined density is bounded near
zero, and $\int |h|\le2$. The pole contribution is $o(1)$: on
$\Im t=\pm1/2$, $|P(t)|\le D^{1/2}\sum|b_R|$, whereas each translated kernel
in $h(\pm i/2)$ has real displacement at least $X-B\ge T$.

On the prime side the total of the off-diagonal terms is
$O(e^{C Y}/B)$ for a fixed $C$. Here are bounds also covering rational
frequencies close to an integer. For $n\le 2D$ and $n\ne R$, the numerator
of $n-R$ in the common denominator is a nonzero integer, hence
$|\log n-\log R|\ge 1/(2D^2)$. The sinc factor costs at most $2D^2/B$,
and $\sum_{n\le2D}(\log n)/\sqrt n\ll D^{1/2}\log(2D)$.
For $n>2D$ the logarithmic difference exceeds $\log2$ and the Fourier
bound gives $O(B^{-1}D^b n^{-b})$; the resulting sum converges since
$b>1/2$. Summing these bounds with $|b_R|$ proves the assertion, uniformly
in the phases and the vector $c$. Choose $\lambda$ so that $C\lambda<1/2$.
The prime diagonal now gives
\[
 2\pi\sum_t z_c(t)P_\pm(t)K_B(t-X)
    =\mp e^{-1/2}|C(v)|V+O(1).
\]
In particular the two positive test functions have signed integrals of
opposite signs and size $\gg\log\log H$.

\subsection{Removing the tails}
The number of zeros of any fixed primitive function in a unit interval
at height $u$, with multiplicity, is $O(\log(2+|u|))$ by the
Riemann--von Mangoldt formula (see also
\href{https://doi.org/10.1090/mcom/3599}{\emph{Counting zeros of Dirichlet
$L$-functions}}). Put $L=H/2-B$. In the $j$th unit interval outside either
end of $(T,T+H]$, the maximal kernel value is
$O(\exp(-c_a e^{a(L+j)}))$, $j\ge0$. This bound also covers a zero
exactly at the excluded endpoint $T$. Summing the unit-interval bounds,
using $H\le T$, gives a tail of size
\[
 O\!\left(\log(2+T)\exp(-c'_a e^{aL})\right)=o(1).
\]
Indeed $e^{aL}\ge(\log\log T)^{1+\eta}$ for some $\eta>0$.
The same estimates apply to negative ordinates. Since $0\le h\le1/B$,
each sign inside the interval has mass $\gg B V$. This proves
Theorem~\ref{th:short}.\hfill$\square$

\paragraph{A single-prime illustration.}
For $\zeta$ and the Dirichlet beta function $L(s,\chi_{-4})$, use the
unnormalized difference of their multiplicities. The unamplified conductor
term is $-\log 4$. Using $P=1\pm\cos(t\log7)$, the two limiting
signed readings after multiplication by $2\pi$ are
$-\log4\mp 2\log7/\sqrt7$. The absolute prime term exceeds $\log4$,
so this one prime already certifies both signs for sufficiently large
windows of the stated length. The growing family of primes in the proof
supplies the factor $\log\log H$.

\section{Opposite multiplicity excesses in every prescribed phase box}
\label{sec:localized}

Fix $q\ge1$ and an odd prime $d\nmid q$. Put $h=\varphi(q)$ and
$\kappa=(\cdot/d)$. For each character $\chi\bmod q$ use its primitive
inducing character $\chi^*$, of conductor $f_\chi$. The character
$\chi^*\kappa$ is primitive of conductor $df_\chi$. At positive height put
\[
 z(\gamma)=\sum_{\chi\bmod q}
 \left(\operatorname{ord}_{1/2+i\gamma}L(s,\chi^*)
      -\operatorname{ord}_{1/2+i\gamma}L(s,\chi^*\kappa)\right),
 \qquad z_\pm=\max(\pm z,0).
\]
Omitted Euler factors have their zeros on $\Re s=0$. Thus the same integer
is obtained with the imprimitive functions. In the range of
\texttt{isPolar\_primeRemainder} this is precisely its residue numerator.
Write $\mathbb T=\mathbb R/(2\pi\mathbb Z)$. For a finite set $B$ of
primes put $\theta_B(t)=(t\log b\bmod 2\pi)_{b\in B}$.

\paragraph{Theorem.}
Let $U$ be a nonempty open box of $\mathbb T^B$, and fix $0<\alpha\le1$.
With Dirichlet GRH as established in the source, there are positive
constants $c,T_0$ depending only on $q,d,B,U,\alpha$ such that, simultaneously
for both signs, $T\ge T_0$ and $T^\alpha\le H\le T$,
\begin{equation}
 \sum_{\substack{T<\gamma\le T+H\\\theta_B(\gamma)\in U}}z_\pm(\gamma)
    \ \ge\ c H\log\log T .
 \label{eq:phase-excess}
\end{equation}
Each distinct ordinate is summed once. The left side counts excess
multiplicity after all coincidences, including coincidences across
different characters, have been cancelled. The case $B=\varnothing$
gives the unrestricted statement. For $q=1$ these are respectively the
zeros and poles, with multiplicities, of $\zeta(s)/L(s,\kappa)$.

\paragraph{A uniform rational-frequency estimate.}
For a primitive character $\xi$ of fixed conductor (including the
zeta case), put
$D_\xi(x)=\sum_{T<\gamma_\xi\le T+H}x^{i\gamma_\xi}$.
If $x=a/b>1$ is in lowest terms with $a,b\le Q$, $Q\ge2$, then
\begin{equation}
 D_\xi(x)=-\frac{H}{2\pi}\frac{\Lambda(x)\xi(x)}{\sqrt{x}}
              +O_\xi(Q^4\log^2(2T)).
 \label{eq:rational-landau}
\end{equation}
Here the numerator in the main term is zero unless $x$ is an integral
prime power. The value at $x^{-1}$ is the conjugate. Banks proves a
sharper uniform estimate for $x\ge2$ in Theorem 1.2 of
\emph{Comparing zeros of distinct Dirichlet $L$-functions}
(\href{https://arxiv.org/abs/2302.07073}{arXiv:2302.07073}).
We give the coarser estimate for the full range $x>1$ because some
ratios of products of primes approach~$1$.

Integrate $x^s L'(s,\xi)/L(s,\xi)$, divided by $2\pi i$, on the
rectangle with vertical sides $-1,2$. Move each horizontal side by
at most $1$ to a height at distance $\gg_\xi 1/\log(2T)$ from all
nearby zero ordinates. The standard local zero count
$N_\xi(t+1)-N_\xi(t-1)\ll_\xi\log(2t)$ permits this choice.
The local partial-fraction expansion for $L'/L$ then gives
$O_\xi(\log^2(2T))$ on the two horizontal sides (Banks, Lemma 2.2).
The horizontal integrals cost $O_\xi(Q^2\log^2(2T))$. Restoring the
endpoints costs $O_\xi(\sqrt{x}\log(2T))$ in the sum of $x^\rho$ by
GRH. On the right side expand $L'/L=-\sum_{n\ge2}\Lambda(n)\xi(n)n^{-s}$.
The diagonal $n=x$ gives the displayed main term before division by
$\sqrt{x}$. For $n\ne x$ the oscillatory integral of $(x/n)^{it}$ is
$O(1/|\log(x/n)|)$. For $x/2<n<2x$ we have $|x-n|\ge1/Q$ and
$|\log(x/n)|\gg Q^{-2}$. Elsewhere the reciprocal is bounded. Since
$\sum_{n\ge2}\Lambda(n)/n^2<\infty$, the total off-diagonal cost is
$O(Q^4)$. On the left side the functional equation and Stirling's
formula give a logarithmic term $-\log(f_\xi t/(2\pi))$, a uniformly
$O_\xi(1/t)$ remainder, and a logarithmic derivative at $2-it$.
Integration by parts bounds the logarithmic term by
$O_\xi(\log(2T)/\log x)$, with $1/\log x\ll Q$. The absolutely
convergent Dirichlet series at $2-it$ has frequencies $\log(nx)\ge
\log 2$, so costs $O(1)$. These estimates prove
\eqref{eq:rational-landau} after division by $\sqrt{x}$. At frequency
$1$ use the Riemann--von Mangoldt formula. Subtracting the equal numbers
of degree-one functions cancels the common $t\log t$ term and gives
\begin{equation}
  \sum_{T<\gamma\le T+H} z(\gamma)
    =-\frac{hH}{2\pi}\log d+O_{q,d}(\log(2T)).
 \label{eq:conductordifference}
\end{equation}

\paragraph{Choosing a bounded test with a growing prime contribution.}
Let $Y=c_0\log T$ for a sufficiently small fixed $c_0>0$ to be specified
by the error bound below. Let $S$ consist of primes $p\le Y$ outside
$B$ satisfying $p\equiv1\pmod q$ and $\kappa(p)=-1$.
Define $w_p=(\log p)/\sqrt p$, $\sigma^2=\sum_{p\in S}w_p^2$.
This union of reduced classes modulo $qd$ has density $1/(2h)$ among
primes. The prime number theorem in fixed progressions and partial
summation give
\[
 \sigma^2\sim\frac{(\log Y)^2}{4h}.
\]
For large $T$, define the real trigonometric polynomial
\[
 A(t)=\Im\prod_{p\in S}
       \left(1+i\frac{w_p}{\sigma}\cos(t\log p)\right).
\]
Pointwise $|A|\le\exp(\sum_p w_p^2/(2\sigma^2))=\sqrt e$.
Its constant term is zero. The coefficient of
$\cos(t\log p)$ arising from the singleton $p$ is $w_p/\sigma$.
All remaining nonzero terms involve at least two distinct primes
(in fact at least three), each with exponent $1$ or $-1$ in the
frequency ratio. Hence neither this ratio nor its inverse is an integral
prime power. At a singleton $p$ we have
$\sum_\chi\chi^*(p)(1-\kappa(p))=2h$.
Taking the real part in \eqref{eq:rational-landau} consequently gives the
weighted singleton contribution $-hH\sigma/\pi$.

\paragraph{Restricting the phases.}
Take a smooth function $0\le F\le1$ supported inside $U$ with normalized
Haar integral $m>0$ (for empty $B$ take $F=1$). Write $F_M$ for its
rectangular Fourier truncation to exponents $|k_b|\le M$, where
$M=\lfloor\sqrt{\log T}\rfloor$. Smoothness gives
$\|F-F_M\|_\infty\ll_{F,D}M^{-D}$ for every fixed $D$ and a uniformly
bounded sum of absolute Fourier coefficients. The sum of $|z|$ in the
window is $O_{q,d}(H\log T+\log T)$ by the zero count. Thus $F_M$ can
replace $F$ in sums against $z$ and $zA$ with an error $o(H)$.
Here and below $F$ in a sum over ordinates stands for $F(\theta_B(\gamma))$.

In $F_M A$ every frequency with a nonzero exponent in $B$ also has
a nonzero exponent in $S$. Unique factorization makes each such
frequency, and its reciprocal, ineligible for a prime-power main term.
Only the constant Fourier coefficient $m$ of $F_M$ couples to a
singleton in $A$. For $F_M$ alone the constant frequency contributes
$O_{q,d}(H+\log T)$ by \eqref{eq:conductordifference}. Its nonconstant
prime-power terms contribute $O_{F,q,d,B}(H)$, since the factor
$\log b/b^{|k|/2}$ is bounded for the fixed primes $b\in B$ and the
Fourier coefficients are absolutely summable.

For completeness, all the rational frequencies occurring here, in
lowest terms, have numerator and denominator at most
\[
 Q=\exp\!\left(\sum_{p\le Y}\log p+M\sum_{b\in B}\log b\right)
       \le \exp(2Y+O_B(M))
\]
for large $T$. The sum of absolute Fourier coefficients of $A$ is at
most $\prod_{p\in S}(1+w_p/\sigma)\le\exp(\sqrt{|S|})$ by Cauchy--Schwarz.
Using \eqref{eq:rational-landau} termwise therefore costs
$O_{F,q,d}(Q^4\log^2(2T)\exp(\sqrt{|S|}))$. Choose, for example,
$c_0=\alpha/32$. This error is at most $T^{\alpha/4+o(1)}=o(H)$.
We have proved the two estimates
\[
 \sum_{T<\gamma\le T+H}z(\gamma)F=O(H),\qquad
 \sum_{T<\gamma\le T+H}z(\gamma)F A(\gamma)
        =-\frac{mh}{\pi}H\sigma+o(H).
\]
The functions $F(\sqrt e-A)$ and $F(\sqrt e+A)$ take their values
in $[0,2\sqrt e]$ and vanish outside the box. Their signed sums are,
respectively, $+(mh/\pi)H\sigma+O(H)$ and
$-(mh/\pi)H\sigma+O(H)$. Dividing the first by $2\sqrt e$ bounds the
positive mass from below, and negating the second does the same for the
negative mass. Since $\sigma\asymp\log\log T$,
\eqref{eq:phase-excess} follows.

\paragraph{Relation to earlier work.}
The signed list $z(\gamma)/h$ already occurs in the ProveIt archive
\texttt{grh-torsion-and-zero-data.zip}, where it is used to recover
$q$ and $d$. The present question asks where its two signs persist.
Liu and Zaharescu study smooth tests on multidimensional tori in races
between the zeros of two Dirichlet functions
(\href{https://doi.org/10.1016/j.jmaa.2020.124591}{JMAA 494 (2021), 124591}).
Banks compares zero multisets in short height intervals.
Bombieri and Perelli's \emph{Distinct zeros of $L$-functions}
(\href{https://doi.org/10.4064/aa-83-3-271-281}{Acta Arith. 83 (1998), 271--281})
treats large global counts for pairs of $L$-functions under specified
hypotheses. Here the cancellation is between two whole products,
the angular box is fixed arbitrarily before the amplifier is chosen,
and the same height window works in both directions.
This proof makes no priority claim. Its count allows different
multiplicities at a common zero. Simplicity, disjointness of zero sets,
and the conjectural $H\log T$ scale for almost all excess multiplicity
are stronger assertions.

\section{An exact choice of weights for Gaussian Hecke characters}
Let $K=\mathbb Q(i)$. Every nonzero ideal has a generator $a$, and
$\eta((a))=(a/|a|)^4$ is independent of the choice of generator.
Write $L_k(s)=L_K(s,\eta^k)$ and $L_0=\zeta_K$.
For $k\ne0$ these are entire Hecke functions of nonzero infinity type.
Conjugation of ideals gives $L_{-k}=L_k$.
Their zeros enter, for example, the explicit formula for counts of Gaussian prime angles.
See \href{https://arxiv.org/abs/1705.07498}{Rudnick and Waxman, \emph{Angles of Gaussian primes}}.

For a real even trigonometric polynomial on the unit circle write
$P(z)=h_0+\sum_{k=1}^N h_k(z^k+z^{-k})$.
In the half-plane $\Re s>1$, termwise logarithmic differentiation gives
\[
 D_P(s)=\sum_{\mathfrak p,m\ge1}
 \frac{\log N\mathfrak p\,P(\eta(\mathfrak p)^m)}{(N\mathfrak p)^{ms}}
 =-h_0\frac{\zeta'_K}{\zeta_K}(s)
 -2\sum_{k=1}^N h_k\frac{L'_k}{L_k}(s).
\]
Here $\mathfrak p$ ranges over prime ideals, including the prime above 2. A nonnegative $P$
gives nonnegative local weights. Nonnegative $h_k$ ensure that orders of zeros
add with the same sign. Its mean $h_0$ is the residue of $D_P$ at 1.
Thus one can ask for the smallest main term compatible with prescribed
sensitivity to several $L_k$ at once.

\subsection{An optimal polynomial for a block of modes}
Fix $0<\alpha<1$ and integers $1\le r\le N$. Define rising factorials
$(a)_0=1$, $(a)_j=a(a+1)\cdots(a+j-1)$ and set
\[
 c_j=\frac{(\alpha)_j}{j!},\qquad b_j=\frac{(2\alpha)_j}{j!},\qquad
 d=\min\{r,\lfloor N/2\rfloor+1\}.
\]
\noindent\textbf{Theorem.}
Among nonnegative $P$ of degree at most $N$ satisfying $h_{N-j}\ge b_j$ for
$0\le j<r$, the exact minimum of $h_0$ is
\[
 M_{N,r}(\alpha)=2\sum_{j=0}^{d-1}c_j^2
 -\boldsymbol{1}_{\{N=2d-2\}}c_{d-1}^2.
\]
It is attained with every $h_k\ge0$. An extremizer is
$P(z)=|\sum_{i=0}^N v_i z^i|^2$, where $v_i=c_{\min(i,N-i)}$ if
$\min(i,N-i)<d$ and $v_i=0$ otherwise.

The parameter $\alpha=1/2$ gives equal required gains $b_j=1$. With $r=N$
this is precisely Belov's theorem after multiplying his cosine polynomial by two.
Belov also proved uniqueness for that case. The weighted block and its use in
the Hecke criterion are developed here. No priority claim is made for the general formula.
See \href{https://doi.org/10.1070/IM1994v043n03ABEH001582}
{Belov, \emph{On an extremal problem on the minimum of a trigonometric polynomial}}.

\medskip\noindent\emph{Proof.}
The coefficients $c_i$ decrease and $(1-z)^{-\alpha}=\sum c_i z^i$,
so $\sum_{i=0}^j c_ic_{j-i}=b_j$. By symmetry,
$h_{N-j}=\sum_{i=0}^j v_i v_{j-i}$.
For $j<d$ this equals $b_j$. If $d\le j<r$, the support covers every index
and $v_i\ge c_i$ throughout. The remaining constraints follow. Parseval gives the stated
mean and every autocorrelation $h_k$ is nonnegative.

For the lower bound use the Fej\'er--Riesz factorization
$P(z)=|\sum_{i=0}^N u_i z^i|^2$, padding with zeros when needed.
Put $A_i=|u_i|$ and $B_i=|u_{N-i}|$ for $0\le i<d$.
Set $w_i=1$, except that $w_{d-1}=1/2$ when $N=2d-2$.
Then $\sum w_i(A_i^2+B_i^2)\le h_0$, since the center index counts once.
Write $e_h=[z^h](1-z)^\alpha$, set $y_i=w_ic_i$, and define
\[
 \lambda_j=\sum_{h=0}^{d-1-j}e_h y_{j+h},\qquad
 H_{ij}=\begin{cases}\lambda_{i+j}& i+j<d,\\0&i+j\ge d.\end{cases}
\]
Here $e_0=1$, $e_h<0$ for $h\ge1$, and $\sum_{h=0}^m e_h>0$,
as follows by taking the coefficient at $z^m$ of $(1-z)^{\alpha-1}$.
Since $y_i$ is positive and nonincreasing, $\lambda_j>0$. The convolution
of $c$ with $e$ is $(1,0,\ldots)$, giving
$(Hc)_i=\sum_{j<d-i}c_j\lambda_{i+j}=w_ic_i$.

For any $i,j$, the arithmetic--geometric mean inequality gives
$2 A_i B_j\le A_i^2 c_j/c_i+B_j^2 c_i/c_j$. Multiply by $H_{ij}$
and sum, using symmetry and $Hc=(w_i c_i)_i$. It follows that
\[
 2\sum_{j<d}\lambda_j b_j
 \le 2\sum_{j<d}\lambda_j h_{N-j}
 \le 2\sum_{i,j<d} H_{ij}A_i B_j
 \le \sum_{i<d}w_i(A_i^2+B_i^2)\le h_0.
\]
The first sum on the left, without the factor two, is
$\sum_{i,j<d}H_{ij}c_i c_j=\sum_{i<d}w_ic_i^2$.
This proves sharpness.\hfill$\square$

\medskip
For example, for $N=r=4$ and $\alpha=1/3$ the factor is
$1+z/3+2z^2/9+z^3/3+z^4$, and the optimal mean is $184/81$.
It meets the first three upper-frequency constraints with equality.
The lower-frequency constraints beyond $\lfloor N/2\rfloor+1$ never increase the optimum.
Stirling's formula gives $c_j\sim j^{\alpha-1}/\Gamma(\alpha)$. For the full
block as $N\to\infty$ this yields
\[
 M_{N,N}(\alpha)=
 \begin{cases}
 2\sum_{j\ge0}c_j^2+o(1), &0<\alpha<1/2,\\
 (2/\pi)\log N+O(1), &\alpha=1/2,\\
 \dfrac{2(N/2)^{2\alpha-1}}{(2\alpha-1)\Gamma(\alpha)^2}(1+o(1)),
 &1/2<\alpha<1 .
 \end{cases}
\]
The bounded-cost regime prescribes decreasing gains. Each of the $N$ gains
is strictly positive at every fixed $N$.

\subsection{The arithmetic estimate required by the extremizer}
Set $g_P=\lim_{s\to1}(D_P(s)-h_0/(s-1))-h_0$, and define
\[
 U_P(x)=\sum_{(N\mathfrak p)^m\le x}\!
 \log N\mathfrak p\,P(\eta(\mathfrak p)^m)
 \left(\frac1{(N\mathfrak p)^m}-\frac1x\right)-h_0\log x-g_P.
\]
\noindent\textbf{Criterion.}
Fix the extremal polynomial. An eventual bound $U_P(x)\le C_\epsilon
x^{-1/2+\epsilon}$ for each $0<\epsilon<1/2$ implies RH for
each $L_k$ with $h_k>0$. In particular, $r=N$ detects every
$L_k$ for $1\le k\le N$.

\medskip\noindent\emph{Proof.}
For $\Re s>1$, termwise integration gives
\[
 \int_1^\infty U_P(x)x^{-s}\,dx=
 \frac{D_P(s)}{s(s-1)}-\frac{h_0}{(s-1)^2}-\frac{g_P}{s-1}.
\]
The expression is regular at 1 by construction. At a zero $\rho$
in the open strip, its residue is
$(h_0\operatorname{ord}_\rho\zeta_K+2\sum h_k\operatorname{ord}_\rho L_k)/
[\rho(1-\rho)]$. This is positive for a real zero detected by $P$, and
nonzero at any detected nonreal zero. The Hecke functions have no zeros on
$\Re s=1$.
Set $\beta=1/2+\epsilon$. Above a suitable cutoff the integrand
$C_\epsilon x^{\beta-1}-U_P(x)$ is nonnegative. Its Mellin transform
continues meromorphically into $\Re s>\beta$. A real pole there would
have negative residue. If its abscissa of convergence exceeded $\beta$,
Landau's theorem would give a singularity at that real abscissa, where
a negative residue contradicts nonnegativity when approached from the right.
Thus its integral is holomorphic on $\Re s>\beta$, excluding every detected zero there.
Let $\epsilon$ decrease to zero and apply the Hecke functional equations.\hfill$\square$

The optimum determines the smallest pole term allowed by this class of detector.
The arithmetic estimate still concerns the angular values on individual Gaussian
prime powers. The source's integral theta-lift comparison for a quadratic
character and a congruence class supplies an estimate for different local coefficients.
For the proposed transfer the unresolved task is to realize an angular weight
in that comparison, already for $N=r=1$.

\paragraph{Further questions.}
Does a bounded-length interval suffice for each directional multiplicity excess for a fixed pair? Can one force a genuinely unshared zero in every triple-logarithmic interval? For Gaussian prime angles, can the local divisibility calculation in the theta lift be made sensitive to even one nonconstant angular harmonic? These ask respectively for a stronger localization argument, information on individual orders at common zeros, and an arithmetic realization of the angular weights.
\end{document}
